% !TEX program = lualatex
%!TEX TS-program = lualatex
\documentclass[12pt,a4paper]{article}
\usepackage{fontspec}
\usepackage{polyglossia}
\setdefaultlanguage{polish}
\usepackage{geometry}
\geometry{margin=2.1cm}
\usepackage[table]{xcolor}
\usepackage{longtable}
\usepackage{array}
\usepackage{booktabs}
\usepackage{enumitem}
\usepackage{ragged2e}
\definecolor{tableHeader}{HTML}{E8EEF7}
\definecolor{tableStripe}{HTML}{F7F9FC}
\newcolumntype{L}[1]{>{\RaggedRight\arraybackslash}p{#1}}
\newcolumntype{C}[1]{>{\Centering\arraybackslash}p{#1}}
\renewcommand{\arraystretch}{1.22}
\setlength{\tabcolsep}{4pt}
\setlist[itemize]{leftmargin=*, itemsep=2pt, topsep=2pt}
\setlength{\parindent}{0pt}
\setlength{\parskip}{6pt}
\emergencystretch=3em
\sloppy
\begin{document}

\section*{Raport ze sprintu nr 4 - Calisia Banking App}
Data raportu: 18.05.2026

Okres sprintu: 13.05 - 18.05.2026

Product Owner / Scrum Master: Ksenia Makarevich

\subsection*{Lista osób}
\begin{itemize}
\item Kubryn Jeremiasz
\item Łaszuk Jakub
\item Makarevich Ksenia
\item Nuckowski Michał
\item Ruczyńska Nikola
\item Wołczko Jakub
\item Olszewski Dawid
\end{itemize}

\subsection*{Product Owner i Scrum Master sprintu}
Ksenia Makarevich

\subsection*{Lista historyjek na sprint}
\begingroup
\footnotesize
\rowcolors{2}{tableStripe}{white}
\begin{longtable}{@{}C{0.08\textwidth}L{0.62\textwidth}C{0.10\textwidth}L{0.12\textwidth}@{}}
\toprule
\rowcolor{tableHeader}
\textbf{Kod} & \textbf{Historyjka} & \textbf{SP} & \textbf{Status} \\
\midrule
\endfirsthead
\toprule
\rowcolor{tableHeader}
\textbf{Kod} & \textbf{Historyjka} & \textbf{SP} & \textbf{Status} \\
\midrule
\endhead
S4-01 & Otwarcie rachunku bankowego z nazwą, numerem, walutą i typem konta & 8 & Done \\
S4-02 & Walidacja numeru rachunku, waluty i typu konta & 5 & Done \\
S4-03 & Wypłata środków z rachunku & 8 & Done \\
S4-04 & Blokada wypłaty przekraczającej saldo & 5 & Done \\
S4-05 & Obsługa wpływu przychodzącego na rachunek & 5 & Done \\
S4-06 & Zapisywanie wpisów transakcyjnych dla wpłat i wypłat & 8 & Done \\
S4-07 & Testy rachunków, wpłat i wypłat & 5 & Done \\
\bottomrule
\end{longtable}
\endgroup

\subsection*{Podsumowanie sprintu}
Cel sprintu: Rachunki bankowe i operacje pieniężne. Udostępnienie klientowi podstawowych produktów bankowych oraz operacji zmieniających saldo.

Zakres sprintu obejmował 7 historyjek o łącznej estymacji 44 SP. Wszystkie historyjki zostały dowiezione ze statusem Done, dlatego osiągnięta prędkość sprintu wyniosła 44 SP.

Udostępniono otwarcie rachunku bankowego przez chroniony endpoint API, generowanie numeru produktu, wpływ przychodzący oraz wypłatę środków. System zapisuje saldo i wpisy transakcyjne, a reguły domenowe blokują wypłatę przekraczającą dostępne środki.

Sprint przesunął aplikację z etapu konta klienta do realnych operacji finansowych. Powstał fundament pod dashboard, historię zdarzeń oraz przyszłe przelewy.

\subsubsection*{Najważniejsze rezultaty}
\begin{itemize}
\item Dodano otwieranie rachunku bankowego.
\item Obsłużono wpływ i wypłatę środków.
\item Wdrożono blokadę wypłaty ponad saldo.
\item Zapisano wpisy transakcyjne potrzebne do historii operacji.
\end{itemize}

\subsection*{Retrospekcja sprintu}
Retrospekcja wskazała, że złożoność domeny finansowej wymaga wcześniejszego doprecyzowania nazewnictwa i scenariuszy testowych.

\subsubsection*{Co się udało}
\begin{itemize}
\item Sprawna implementacja otwierania rachunku bankowego zgodnie z CQRS.
\item Dobre wykorzystanie reguł domenowych do kontroli salda.
\item Poprawne rozdzielenie operacji wpłaty i wypłaty.
\item Testy potwierdziły najważniejsze ścieżki rachunków i operacji pieniężnych.
\end{itemize}
\subsubsection*{Poziom energii zespołu}
Średni poziom energii zespołu: 4,0/5.

Energia była dobra i zadaniowa. Zakres był większy niż wcześniej, dlatego wymagał dokładniejszej koordynacji backendu, testów i scenariuszy API.

\subsubsection*{Czego udało się nauczyć}
\begin{itemize}
\item Zespół lepiej zrozumiał modelowanie operacji finansowych w domenie.
\item Większe historyjki wymagają dokładniejszego rozbicia technicznego.
\item Kolekcję Bruno warto aktualizować od razu po zmianach endpointów.
\end{itemize}
\subsubsection*{Wnioski i działania na kolejny sprint}
\begin{itemize}
\item Przygotować mapę endpointów rachunkowych dla testerów.
\item Dopisać w Bruno osobne scenariusze dla rachunku, wpływu i wypłaty.
\item Przed kolejnym sprintem ustalić strukturę modelu odczytowego dashboardu.
\end{itemize}

\subsection*{Następny sprint}
Product Owner / Scrum Master w następnym sprincie: Michał Nuckowski

Główny cel następnego sprintu: Dashboard i model odczytowy. Pokazanie klientowi aktualnego stanu produktów, salda łącznego i ostatnich zdarzeń w panelu bankowości.

\end{document}
